\chapter{Pre-Rendering and Cosmic Assets}

\begin{chaptermapbox}
\textbf{In this chapter:}
\begin{compactchecklist}
\item Asset baking: why game developers do expensive work once, before anyone plays
\item Cosmic inflation as the universe's baking step, and the problems it solves
\item The cosmic microwave background: a nearly uniform texture with faint, pre-set patterns
\item Procedural generation, random seeds, and how a few numbers grow into the cosmic web
\item What inflation explains well, and what remains uncertain
\end{compactchecklist}
\end{chaptermapbox}

When you load a modern video game and walk into a sunlit room, the shadows under the furniture look soft and realistic. Light seems to bounce off the walls and gently brighten the corners. Calculating that kind of lighting properly, tracing how every ray of light reflects around the room, is enormously expensive. Doing it live, sixty times a second, would bring most machines to their knees.

So, for scenery that doesn't move, developers usually don't. They calculate the lighting once, during development, on powerful machines that may run for hours. Then they ``bake'' the result into the textures of the walls and floor. When you play, the game just paints those pre-lit textures onto the surfaces. The expensive physics was done before you arrived. At runtime, all that's left is a cheap lookup.

The same thinking runs through all of game design. Heavy computation is moved to the beginning, done once, and stored as a simple asset that the running game can reuse for free.

This chapter argues that the universe's opening moments look a lot like that. In the first tiny fraction of a second, a process called cosmic inflation seems to have set up starting conditions that were uniform, simple, and governed by just a few numbers. Everything since has been built on those baked-in assets, using a technique game developers would recognize instantly: procedural generation from a seed.

\section{Cosmic Inflation as Asset Baking}

\subsection{Three Puzzles About the Beginning}

By the 1970s, the Big Bang model was well established. The universe had been expanding from a hot, dense state for billions of years. The cosmic microwave background, discovered in 1965, was the leftover glow of that early heat. But the model had three nagging problems.

\textbf{The horizon problem.} The cosmic microwave background has the same temperature in every direction to about one part in a hundred thousand \cite{smoot1992}. That's astonishingly uniform. Two patches on opposite sides of the sky are so far apart that, in the standard Big Bang picture, light had not had time to travel between them by the time that radiation was released. Chapter~5 explained that nothing can share information faster than light. So how did these regions, which were never in contact, end up at exactly the same temperature? It's like finding two strangers on opposite sides of the world who have never communicated, and discovering that they've set their thermostats to the same value to five significant figures.

\textbf{The flatness problem.} The overall geometry of the universe is extremely close to flat, meaning parallel lines stay parallel over cosmic distances. In the standard model, any small departure from flatness in the early universe would have grown enormously over time. For the universe to be as flat as it is today, it must have started out flat to an absurd precision, something like one part in \(10^{60}\).

\textbf{The relic problem.} Many theories of particle physics predict that the hot early universe should have produced heavy, exotic particles, such as magnetic monopoles, in large numbers. We've never found a single one.

\subsection{Guth's Idea}

In 1981, Alan Guth proposed a solution to all three at once \cite{guth1981}. Suppose that in the first instant after the Big Bang, the universe went through a brief period of extraordinarily rapid expansion. In something like \(10^{-32}\) seconds, it grew by a factor of at least \(10^{26}\), from smaller than a proton to something perhaps the size of a grapefruit or larger. Andrei Linde and others refined the idea in the following year into a form that works better in detail \cite{linde1982}.

That expansion solves the puzzles neatly.

\begin{compactchecklist}
\item \textbf{Horizon:} Before inflation, the entire observable universe was a tiny region small enough to come to a uniform temperature. Inflation then stretched it far beyond the distance light could later cross. The distant patches of sky look identical because they were once neighbors.
\item \textbf{Flatness:} Blowing up any curved surface enormously makes it look flat, the way the Earth looks flat from where you stand. Sufficient inflation drives the observable spatial curvature toward zero; the conditions needed to start inflation remain a separate question.
\item \textbf{Relics:} Any exotic particles produced before inflation were diluted so thoroughly that we'd expect to find essentially none in our observable universe.
\end{compactchecklist}

\subsection{Reading It as a Bake Step}

Now look at what inflation accomplishes from the engine's point of view.

Before inflation, there may have been any amount of irregularity, any messy distribution of energy, any number of exotic relics. After inflation, the observable universe is smooth, flat, and uniform, with nearly identical conditions everywhere. The messy, expensive-to-describe initial state has been replaced by a clean, cheap one.

When inflation ended, the energy that had driven it was converted into a hot soup of particles, a process called reheating. That soup quickly reached thermal equilibrium. And as Chapter~6 explained, equilibrium is the ultimate cheap description. A hot plasma in equilibrium is described by its temperature and a few numbers for its composition.

A thermal description of that plasma needs a temperature and composition parameters, not a list of particle positions. That is a saving in our description. It does not show that inflation replaced a microscopic state with a small data file.

\subsection{The First Three Minutes}

The baking metaphor gets sharper when you look at what happened next. Within the first few minutes, as the universe cooled, protons and neutrons fused into the first light nuclei: hydrogen, helium, and traces of lithium. This is called Big Bang nucleosynthesis.

With the expansion model and reaction physics specified, the baryon-to-photon ratio is a central cosmological input to nucleosynthesis. It is about six parts in ten billion; nuclear reaction rates and the neutron lifetime also matter. Calculations predict that roughly a quarter of ordinary matter by mass becomes helium-4, with most of the rest remaining hydrogen \cite{pitrou2018nucleosynthesis}. Deuterium and helium provide successful tests, while inferred primordial lithium-7 is substantially below the standard prediction. That discrepancy prevents an unqualified claim of agreement for all light elements.

A small set of physical inputs constrains the abundance pattern. That is predictive power for statistical quantities, not a replacement for every microscopic reaction.

\subsection{The Cosmic Microwave Background as a Texture}

About 380,000 years after the Big Bang, the universe had cooled enough for electrons and nuclei to combine into neutral atoms. At that moment, light was suddenly free to travel without constantly scattering off charged particles. That light has been traveling ever since. Stretched by the expansion of the universe, it now arrives as microwaves at a temperature of about 2.7 kelvin. That's the cosmic microwave background.

For a game designer, the cosmic microwave background looks like a baked texture wrapped around the whole sky. It's nearly uniform. It's described almost entirely by one number, its temperature. And it was set once, long ago, and its observed pattern also carries later effects, including gravitational lensing and scattering. The texture comparison is therefore limited.

But it isn't perfectly uniform. That's where the story gets interesting.

\section{Procedural Generation}

\subsection{How Games Build Worlds from Seeds}

Some of the largest game worlds ever built were never designed by hand. They were generated.

In a procedurally generated game, the developers don't model each mountain and valley. They write a set of rules for producing terrain, and they feed those rules a random number called a seed. The rules then deterministically expand the seed into a full world: hills, rivers, forests, caves. Use the same seed, and you get exactly the same world. Use a different seed, and you get a different but equally plausible one.

A key ingredient in many of these systems is smooth random noise. In 1985, Ken Perlin introduced a technique, now called Perlin noise, for generating random-looking but natural patterns, with structure at many different scales \cite{perlin1985}. Layer noise at several scales and you get something that looks like a coastline, a cloud, or a mountain range. Perlin later received an Academy Award for technical achievement for this work, because it transformed how films render natural textures.

The economics are extraordinary. A procedurally generated world can be vastly larger than anything that could be stored. The rules and the seed are all you keep. The world itself is computed as needed, which, as Chapter~4 discussed, is the essence of lazy evaluation.

\subsection{The Universe's Seed}

Back to the cosmic microwave background. In 1992, the COBE satellite detected tiny variations in its temperature: some patches slightly warmer, some slightly cooler, by about one part in a hundred thousand \cite{smoot1992}. Since then, the WMAP and Planck satellites have mapped those variations in exquisite detail \cite{planck2018}.

These faint ripples are the seeds of everything. Slightly denser regions in the early universe had slightly stronger gravity. They pulled in surrounding matter, became denser still, and pulled in more. Over hundreds of millions of years, the densest regions collapsed into the first stars and galaxies. Over billions of years, gravity built the vast pattern of filaments, clusters, and empty voids that astronomers call the cosmic web.

Where did the ripples come from? Inflation offers an answer, and it's one of the most beautiful ideas in cosmology. During inflation, tiny quantum fluctuations, the unavoidable jitter that Chapter~4 discussed, were stretched from microscopic to astronomical sizes \cite{mukhanov1981}. Every galaxy, including our own, may have begun as a quantum fluctuation smaller than an atom, inflated to cosmic scale.

\subsection{Six Numbers for the Whole Universe}

The statistical description is short even though a particular realization contains far more detail.

The statistical pattern of the ripples is remarkably simple. To an excellent approximation, it's what statisticians call a Gaussian random field: a random field specified statistically by its correlations. This resembles procedural terrain at the level of generating patterns, not as an identification with a game algorithm. Its character is set by two numbers. One gives the overall strength of the ripples. The other, called the spectral index, describes how that strength changes from large scales to small ones. The Planck satellite measured it at about \(0.965\), very close to but measurably different from \(1\), which is close to what simple models of inflation predict \cite{planck2018}.

Add four more parameters, covering the amounts of ordinary matter and dark matter, the expansion rate, and the effect of the early universe's free electrons on the light passing through, and you have what cosmologists call the six-parameter standard model of cosmology. Those six numbers fit the detailed pattern of the cosmic microwave background, measured at millions of points across the sky, with remarkable accuracy.

\begin{examplebox}
Compare two ways to describe the early universe.

\textbf{Brute force:} The density and velocity of matter at every point in a region now billions of light-years across. Effectively infinite data.

\textbf{Procedural:} Six numbers for the rules, plus one particular random realization of the ripples, the ``seed.'' The rules are tiny. The seed is the only part that carries much information.

A numerical seed can generate a model realization, but no compact seed for the exact observed universe has been established. Simulations reproduce selected statistics to useful accuracy; they do not reproduce every galaxy or every observable exactly.
\end{examplebox}

\subsection{Growing the Cosmic Web}

That last point has been tested directly. In 2005, Volker Springel and colleagues ran the Millennium Simulation \cite{springel2005}. They started ten billion simulated particles in a cube two billion light-years across, with initial ripples matching the statistics measured from the cosmic microwave background. Then they let gravity act for the equivalent of 13 billion years.

What emerged looked like our universe. Filaments of galaxies stretched across the box. Dense clusters formed where filaments met. Huge empty voids opened between them. The statistical pattern matched what galaxy surveys see in the real sky.

No one put the filaments there by hand. They grew from the seed, according to the rules, just as a procedurally generated mountain range grows from its noise function.

\subsection{Stars from a Short Recipe}

Procedural generation continues all the way down. Once gravity collects a cloud of gas, the cloud collapses and forms stars. And for a star, an old rule of thumb in astrophysics says that its entire life story is determined almost entirely by just two things: its starting mass and its chemical composition. Give a stellar astrophysicist those two numbers, and they can tell you how hot the star will burn, how long it will live, and whether it will end as a white dwarf, a neutron star, or a black hole.

It's another instance of the same pattern. Enormous complexity, generated by simple rules from a few input numbers.

\begin{quickpracticebox}
Run two model universes with the same cosmological parameters but different random realizations. They can have similar clustering statistics while putting individual galaxies in different places. Matching the statistical pattern supports the model at that level; it does not recover our universe's exact initial state. Ask which comparison an advertised simulation has actually passed.
\end{quickpracticebox}

\section{What Inflation Does and Doesn't Explain}

The baking-and-seeding picture is appealing, and it's grounded in real, precise measurements. But it would be easy to oversell it.

\begin{researchnotebox}
\textbf{Well supported:} The cosmic microwave background is uniform to about one part in a hundred thousand, its ripples follow nearly Gaussian statistics with a spectral index near \(0.965\), and the six-parameter model fits the data well \cite{planck2018}. Big Bang nucleosynthesis successfully accounts for key deuterium and helium measurements, while the lithium discrepancy remains a limitation \cite{pitrou2018nucleosynthesis}.

\textbf{Promising but limited:} Inflation is the leading explanation for the uniformity, flatness, and ripple pattern, and its predictions have matched observations. Many cosmologists consider it well motivated.

\textbf{Open:} Inflation has not been confirmed by a ``smoking gun'' observation. Many models predict a faint pattern of primordial gravitational waves in the cosmic microwave background. It has not been detected. A joint BICEP2/Keck--Planck analysis found significant polarized dust and no statistically significant evidence for primordial gravitational waves in those data \cite{bicep2015joint}. Some physicists argue inflation is too flexible to be falsified, and alternative models exist.
\end{researchnotebox}

There's also a deeper puzzle that inflation doesn't fully solve. For entropy to have been increasing ever since the beginning, as Chapter~6 described, the universe must have started in a state of extraordinarily low entropy. A smooth, uniform distribution of matter is low entropy in a gravitating universe, because gravity tends to pull matter into clumps, and clumping releases heat and raises entropy. Roger Penrose has long argued that the specialness of this initial state is a profound mystery that inflation pushes back but doesn't answer \cite{penrose1989}.

In engineering terms, inflation explains the baking step. It doesn't fully explain why the oven was set up so neatly in the first place. Philosophers sometimes call the assumption of a low-entropy beginning the ``Past Hypothesis'' \cite{albert2000}, and Chapter~10 returns to it, because it's the key to the arrow of time.

\section{A Universe That Doesn't Need a Full Save File}

The distinction to retain is between simple statistics and a complete initial state.

Inflation offers a mechanism for the observed large-scale smoothness and the statistics of primordial fluctuations. Those statistics can be described compactly. The particular realization, and the conditions required for inflation to begin, carry questions that the short parameter list does not answer.

The baking analogy describes a proposed early physical process. It provides no evidence that the universe stored a reusable asset or avoided computing its subsequent evolution.

That handles the beginning. The final chapter of Part V turns to the long run. A universe running for billions of years, and possibly forever, faces two remaining problems: how to keep its far-flung regions from becoming one enormous, interconnected calculation, and how to stop its history from piling up without limit.

\begin{keytakeawaybox}
A few parameters describe important cosmic statistics, not the exact cosmic state; inflation is a supported proposal for their origin, and asset baking remains an analogy.
\end{keytakeawaybox}
